\documentclass[twoside,openright,titlepage,fleqn,
	headinclude,11pt,a4paper,BCOR5mm,footinclude
	]{scrbook}
%--------------------------------------------------------------
        \newcommand{\myTitle}{Progettazione e Analisi di Algoritmi\xspace}
% use the right myDegree option
\newcommand{\myDegree}{Corso di Laurea Magistrale in Informatica\xspace}
%\newcommand{\myDegree}{
	%Corso di Laurea Specialistica in Scienze e Tecnologie 
	%dell'Informazione\xspace}
\newcommand{\myName}{Massimo Nocentini\xspace}
\newcommand{\myProf}{Donatella Merlini\xspace}
\newcommand{\myOtherProf}{Maria Cecilia Verri\xspace}
\newcommand{\mySupervisor}{Nome Cognome\xspace}
\newcommand{\myFaculty}{
	Facolt\`a di Scienze Matematiche, Fisiche e Naturali\xspace}
\newcommand{\myDepartment}{
	Dipartimento di Sistemi e Informatica\xspace}
\newcommand{\myUni}{\protect{
	Universit\`a degli Studi di Firenze}\xspace}
\newcommand{\myLocation}{Firenze\xspace}
\newcommand{\myTime}{Anno Accademico 2012-2013\xspace}
\newcommand{\myVersion}{Version 0.1\xspace}

%--------------------------------------------------------------
\usepackage[latin1]{inputenc} 
\usepackage[T1]{fontenc} 
\usepackage[square,numbers]{natbib} 
\usepackage[fleqn]{amsmath}  
\usepackage[italian]{babel}
\usepackage{ae,aecompl}
\usepackage{graphicx}
\usepackage{latexsym}
\usepackage{amsmath, amsthm, amssymb}
\usepackage{rotating}
\usepackage{boxedminipage}
\usepackage{multicol}

%--------------------------------------------------------------
\usepackage{dia-classicthesis-ldpkg} 
%--------------------------------------------------------------
% Options for classicthesis.sty:
% tocaligned eulerchapternumbers drafting linedheaders 
% listsseparated subfig nochapters beramono eulermath parts 
% minionpro pdfspacing
\usepackage[eulerchapternumbers,subfig,beramono,eulermath,
	parts]{classicthesis}
%--------------------------------------------------------------
\newlength{\abcd} % for ab..z string length calculation
% how all the floats will be aligned
\newcommand{\myfloatalign}{\centering} 
\setlength{\extrarowheight}{3pt} % increase table row height
\captionsetup{format=hang,font=small}
%--------------------------------------------------------------
% Layout setting
%--------------------------------------------------------------
\usepackage{geometry}
\geometry{
	a4paper,
	ignoremp,
	bindingoffset = 1cm, 
	textwidth     = 13.5cm,
	textheight    = 21.5cm,
	lmargin       = 3.5cm, % left margin
	tmargin       = 4cm    % top margin 
}
%--------------------------------------------------------------
\usepackage{listings}
\usepackage{hyperref}
% My Theorem
\newtheorem{oss}{Observation}[section]
\newtheorem{exercise}{Exercise}[section]
\newtheorem{thm}{Theorem}[section]
\newtheorem{cor}[thm]{Corollary}

\newtheorem{lem}[thm]{Lemma}

\newcommand{\vect}[1]{\boldsymbol{#1}}

% questo comando e' relativo alle correzioni che puo
% apportare il prof se lo desidera.
\newcommand{\prof}[1]{\boldsymbol{#1}}

% instead of boldsymbol I can use the arrow above the letter with
%\newcommand{\vect}[1]{\vec{#1}}

% page settings
% \pagestyle{headings}
%--------------------------------------------------------------
\begin{document}
\frenchspacing
\raggedbottom
\pagenumbering{roman}
\pagestyle{plain}
%--------------------------------------------------------------
% Frontmatter
%--------------------------------------------------------------
%--------------------------------------------------------------
% titlepage.tex (use thesis.tex as main file)
%--------------------------------------------------------------
\begin{titlepage}
	\begin{center}
   	\large
      \hfill
      \vfill
      \begingroup
			\spacedallcaps{\myUni} \\ 
			\myFaculty \\
			\myDegree \\ 
			\vspace{0.5cm}
         \includegraphics[scale=.065]{logo/unifi}\\
         \vspace{0.5cm}    
         %% -------put here the type of document---ie: Tesi di Laurea
         Elaborato d'Esame
      \endgroup 
      \vfill 
      \begingroup
      	\color{Maroon}\spacedallcaps{\myTitle} \\ \bigskip
      \endgroup
      \spacedlowsmallcaps{\myName}
      \vfill  
      % Professoresse:
      \itshape{\myProf}\\\itshape{\myOtherProf}
      \vfill  
      \myTime
      \vfill                      
	\end{center}        
\end{titlepage}   
%--------------------------------------------------------------
% back titlepage
%--------------------------------------------------------------
   \newpage
	\thispagestyle{empty}
	\hfill
	\vfill
	\noindent\myName: 
	\textit{\myTitle,} 
	\myDegree, \textcopyright\ \myTime
%--------------------------------------------------------------
% back titlepage end
%--------------------------------------------------------------
\pagestyle{scrheadings}
%--------------------------------------------------------------
% Mainmatter
%--------------------------------------------------------------
\pagenumbering{arabic}

% settings for the lstlisting environment
\lstset{
	language = R
	, numbers = left 
	, basicstyle=\sffamily%\footnotesize
%	, frame=single
	, tabsize=2
	, captionpos=b
	, breaklines=true
	, showspaces=false
	, showstringspaces=false
}

\tableofcontents

\newpage

%\section*{Introduzione}
Queste note contengono tutto il mio materiale di studio per l'esame di
\emph{Progettazione e Analisi di Algoritmi}.

% \subsection*{Sintassi esercizi}
% Per gli esercizi utilizziamo la sintassi \textbf{Exercise n[(m)]},
% dove $n$ rappresenta la numerazione locale all'interno delle sezioni
% di questo documento, $m$ rappresenta l'identificativo dell'esercizio
% fissato nel documento \\
% \href{http://www.ds.unifi.it/gmm/resources/capitoli.pdf}{
%   http://www.ds.unifi.it/gmm/resources/capitoli.pdf}\\
% Il reference $(m)$ non \`e sempre presente, come evidenziato dall'uso
% delle parentesi quadre.


\newpage

\subsection*{Text contents}
All the text content is distributed under:\\
\textbf{ This work is licensed under the Creative Commons Attribution,
  NonCommercial, ShareAlike 3.0 Unported License. To view a
  copy of this license, visit\\
  \href{http://creativecommons.org/licenses/by-nc-sa/3.0/}{http://creativecommons.org/licenses/by-nc-sa/3.0/}\\
  or send a letter to Creative Commons, 444 Castro Street, Suite 900,
  Mountain View, California, 94041, USA.}

\begin{center}
\includegraphics{cc-icons-eps/by}
\includegraphics{cc-icons-eps/cc}
\includegraphics{cc-icons-eps/nc}
\includegraphics{cc-icons-eps/sa}
\end{center}

\subsection*{Sources}
All sources are distributed under, where the word ``Software'' is
referred to all of the sources that are present in this work: \\
\textbf{
  Copyright (c) 2011 Massimo Nocentini\\
  Permission is hereby granted, free of charge, to any person
  obtaining a copy of this software and associated documentation files
  (the "Software"), to deal in the Software without restriction,
  including without limitation the rights to use, copy, modify, merge,
  publish, distribute, sublicense, and/or sell copies of the Software,
  and to permit persons to whom the Software is furnished
  to do so, subject to the following conditions:\\
  The above copyright notice and this permission notice shall be
  included in all
  copies or substantial portions of the Software.\\
  THE SOFTWARE IS PROVIDED "AS IS", WITHOUT WARRANTY OF ANY KIND,
  EXPRESS OR IMPLIED, INCLUDING BUT NOT LIMITED TO THE WARRANTIES OF
  MERCHANTABILITY, FITNESS FOR A PARTICULAR PURPOSE AND
  NONINFRINGEMENT. IN NO EVENT SHALL THE AUTHORS OR COPYRIGHT HOLDERS
  BE LIABLE FOR ANY CLAIM, DAMAGES OR OTHER LIABILITY, WHETHER IN AN
  ACTION OF CONTRACT, TORT OR OTHERWISE, ARISING FROM, OUT OF OR IN
  CONNECTION WITH THE SOFTWARE OR THE USE OR OTHER DEALINGS IN THE
  SOFTWARE.  }


\chapter{Stima, verifica d'ipotesi e distribuzioni campionarie}



\end{document} 
